% Original PDF page 3; hash in recovery-map.json.
\recoverypage{3}
\section*{3  Compilers, decompilers, and reconstruction loops}
\subsection*{3.1  The deterministic source-to-symbolic path}
\noindent
The local modal codec uses deterministic registry cues, spaCy-derived features, BM25 frame ground-\linebreak[4]
ing, and frame-logic consistency checks. It produces a structured modal declaration before optional\linebreak[4]
expensive model repair. A separate selected structured-text profile has a narrower seven-field rule:\linebreak[4]
modality, actor, action, object, conditions, exceptions, and temporal qualifiers. Its modality vo-\linebreak[4]
cabulary is obligation, permission, and prohibition; its atom vocabulary is supplied by the caller.\linebreak[4]
Unsupported semantics trigger abstention or disclosed partial output. The selected profile explicitly\linebreak[4]
disables learned guidance and repair: training elsewhere in the repository does not imply that this\linebreak[4]
frozen profile silently uses it. [artifact S07,S08,S23,S25]\linebreak[4]
Formalization adapters produce named formula views and source maps. The maps bind IR constituents\linebreak[4]
to source evidence; they are not proof of a correct interpretation. Source-code extraction offers\linebreak[4]
additional candidate security facts, with explicit heuristic and coverage limitations. Media requires a\linebreak[4]
separately measured extraction layer before this logic pipeline; no end-to-end audio/image theorem\linebreak[4]
performance is established here. [artifact S09,S13--S16]\par

\subsection*{3.2  Two different meanings of reconstruction}
\noindent
Symbolic reconstruction follows\par

% Newly transcribed from original PDF equation (1); see recovery-map.json.
\begin{equation}
I_1=C_\psi(x),\qquad \hat{x}=D_\omega(I_1),\qquad I_2=C_\psi(\hat{x}).\tag{1}\label{eq:original-1}
\end{equation}

\noindent
A decompiler renders roles, predicates, modality, conditions, exceptions, and temporal structure\linebreak[4]
into inspectable text. The canonical source-withheld evaluation gives the realizer the IR, not the\linebreak[4]
originating source, its locator, source map, private parser state, or gold target. Comparing I1 and I2\linebreak[4]
tests reconstructable facets and compiler/realizer consistency. It does not prove source fidelity: both\linebreak[4]
programs could consistently omit the same exception. [artifact S07]\linebreak[4]
The adaptive autoencoder's numerical reconstruction is different. Its encode method builds a state\linebreak[4]
containing family predictions and an embedding projection; its decode returns that projection as\linebreak[4]
a numeric vector. Its basic reconstruction target is the sample embedding, not arbitrary source\linebreak[4]
prose. Textual decompilation and structural cycle metrics enter through the codec and multi-view\linebreak[4]
bridge. Calling this vector reconstruction a trained free-text decoder would misstate the inspected\linebreak[4]
implementation. [artifact S24]\par

\subsection*{3.3  Constrained candidates and transfer beyond the legal teacher}
\noindent
The constrained decoder gives learned heads a closed production space: valid-token, grammar, binder,\linebreak[4]
type, and family masks precede candidate scoring. The inspected profile declares a frozen 143-token\linebreak[4]
vocabulary and bounded search, but these declarations are not a learned tokenizer checkpoint or\linebreak[4]
evidence of end-to-end sequence training. The learner may rank admitted productions; a deterministic\linebreak[4]
grammar decides what can be represented. [artifact S33]\linebreak[4]
The original sample builder is legal-specific. A domain-neutral formalization/advisor layer permits\linebreak[4]
other domain adapters to share feature, view, checkpoint, and result contracts, but it does not\linebreak[4]
automatically turn a legal-trained model into a calibrated Security or Intent formalizer. Cross-domain\linebreak[4]
use needs explicit adapter mappings and separate evaluation. Shared infrastructure, statistical transfer,\linebreak[4]
and semantic proof transfer are three distinct claims. [artifact S01,S17,S22]\par

\section*{4  Adaptive multi-view learning}
\subsection*{4.1  Shared parameter state versus per-example memory}
\noindent
The adaptive implementation has sparse parameter groups for feature embeddings, modal families,\linebreak[4]
semantic slots, predicate arguments, logic signatures, view logits, decompiler plans, and family--slot--\linebreak[4]
view interactions. It also retains compatibility fields indexed by sample identity. The generalizable\linebreak[4]
training route disables per-sample memory for updates and evaluation; accepted improvements must\par

